%% ma: L10-B12-0011
%% lop: 10
%% chuong: 5
%% bai: 12
%% dang: 10.12.4
%% loai: TL
%% muc: TH
%% dap_an: "-"
%% nguon: "chuyen-de-cac-so-dac-trung-cua-mau-so-lieu-khong-ghep-nhom-toan-10-chuong-trinh-moi.pdf (GV. Phan Nhật Linh) (P01), Bài 01 Bài tập tự luận 11"
%% kien_thuc: "Lớp 10 Bài 12 (số gần đúng, sai số tuyệt đối, độ chính xác, sai số tương đối, quy tròn số)"
%% trang_thai: chinh-thuc
%% ghi_chu: "Kết luận gốc 'phép đo thứ hai chính xác hơn' sai; đúng là phép đo thứ nhất. Tự duyệt theo yêu cầu giáo viên 22/10/2026."
\begin{ex}
Các nhà thiên văn tính được thời gian để Trái Đất quay một vòng quanh Mặt Trời là $365\pm\dfrac14$ ngày, thời gian để Mặt Trăng quay một vòng quanh Trái Đất là $27{,}3\pm\dfrac1{50}$ ngày. Trong hai phép đo trên, phép đo nào chính xác hơn?
\loigiai{Phép đo thứ nhất: $\delta_1\le\dfrac{0{,}25}{365}\approx0{,}0685\%$.

Phép đo thứ hai: $\delta_2\le\dfrac{0{,}02}{27{,}3}\approx0{,}0733\%$.

Vì $\delta_1<\delta_2$ nên phép đo thứ nhất (Trái Đất quay quanh Mặt Trời) chính xác hơn.}
\end{ex}
